% Options for packages loaded elsewhere
\PassOptionsToPackage{unicode}{hyperref}
\PassOptionsToPackage{hyphens}{url}
%
\documentclass[
  11pt,
  UTF8]{ctexart}
\usepackage{amsmath,amssymb}
\usepackage{iftex}
\ifPDFTeX
  \usepackage[T1]{fontenc}
  \usepackage[utf8]{inputenc}
  \usepackage{textcomp} % provide euro and other symbols
\else % if luatex or xetex
  \usepackage{unicode-math} % this also loads fontspec
  \defaultfontfeatures{Scale=MatchLowercase}
  \defaultfontfeatures[\rmfamily]{Ligatures=TeX,Scale=1}
\fi
\usepackage{lmodern}
\ifPDFTeX\else
  % xetex/luatex font selection
\fi
% Use upquote if available, for straight quotes in verbatim environments
\IfFileExists{upquote.sty}{\usepackage{upquote}}{}
\IfFileExists{microtype.sty}{% use microtype if available
  \usepackage[]{microtype}
  \UseMicrotypeSet[protrusion]{basicmath} % disable protrusion for tt fonts
}{}
\makeatletter
\@ifundefined{KOMAClassName}{% if non-KOMA class
  \IfFileExists{parskip.sty}{%
    \usepackage{parskip}
  }{% else
    \setlength{\parindent}{0pt}
    \setlength{\parskip}{6pt plus 2pt minus 1pt}}
}{% if KOMA class
  \KOMAoptions{parskip=half}}
\makeatother
\usepackage{xcolor}
\usepackage[margin=2.1cm]{geometry}
\usepackage{color}
\usepackage{fancyvrb}
\newcommand{\VerbBar}{|}
\newcommand{\VERB}{\Verb[commandchars=\\\{\}]}
\DefineVerbatimEnvironment{Highlighting}{Verbatim}{commandchars=\\\{\}}
% Add ',fontsize=\small' for more characters per line
\newenvironment{Shaded}{}{}
\newcommand{\AlertTok}[1]{\textcolor[rgb]{1.00,0.00,0.00}{\textbf{#1}}}
\newcommand{\AnnotationTok}[1]{\textcolor[rgb]{0.38,0.63,0.69}{\textbf{\textit{#1}}}}
\newcommand{\AttributeTok}[1]{\textcolor[rgb]{0.49,0.56,0.16}{#1}}
\newcommand{\BaseNTok}[1]{\textcolor[rgb]{0.25,0.63,0.44}{#1}}
\newcommand{\BuiltInTok}[1]{\textcolor[rgb]{0.00,0.50,0.00}{#1}}
\newcommand{\CharTok}[1]{\textcolor[rgb]{0.25,0.44,0.63}{#1}}
\newcommand{\CommentTok}[1]{\textcolor[rgb]{0.38,0.63,0.69}{\textit{#1}}}
\newcommand{\CommentVarTok}[1]{\textcolor[rgb]{0.38,0.63,0.69}{\textbf{\textit{#1}}}}
\newcommand{\ConstantTok}[1]{\textcolor[rgb]{0.53,0.00,0.00}{#1}}
\newcommand{\ControlFlowTok}[1]{\textcolor[rgb]{0.00,0.44,0.13}{\textbf{#1}}}
\newcommand{\DataTypeTok}[1]{\textcolor[rgb]{0.56,0.13,0.00}{#1}}
\newcommand{\DecValTok}[1]{\textcolor[rgb]{0.25,0.63,0.44}{#1}}
\newcommand{\DocumentationTok}[1]{\textcolor[rgb]{0.73,0.13,0.13}{\textit{#1}}}
\newcommand{\ErrorTok}[1]{\textcolor[rgb]{1.00,0.00,0.00}{\textbf{#1}}}
\newcommand{\ExtensionTok}[1]{#1}
\newcommand{\FloatTok}[1]{\textcolor[rgb]{0.25,0.63,0.44}{#1}}
\newcommand{\FunctionTok}[1]{\textcolor[rgb]{0.02,0.16,0.49}{#1}}
\newcommand{\ImportTok}[1]{\textcolor[rgb]{0.00,0.50,0.00}{\textbf{#1}}}
\newcommand{\InformationTok}[1]{\textcolor[rgb]{0.38,0.63,0.69}{\textbf{\textit{#1}}}}
\newcommand{\KeywordTok}[1]{\textcolor[rgb]{0.00,0.44,0.13}{\textbf{#1}}}
\newcommand{\NormalTok}[1]{#1}
\newcommand{\OperatorTok}[1]{\textcolor[rgb]{0.40,0.40,0.40}{#1}}
\newcommand{\OtherTok}[1]{\textcolor[rgb]{0.00,0.44,0.13}{#1}}
\newcommand{\PreprocessorTok}[1]{\textcolor[rgb]{0.74,0.48,0.00}{#1}}
\newcommand{\RegionMarkerTok}[1]{#1}
\newcommand{\SpecialCharTok}[1]{\textcolor[rgb]{0.25,0.44,0.63}{#1}}
\newcommand{\SpecialStringTok}[1]{\textcolor[rgb]{0.73,0.40,0.53}{#1}}
\newcommand{\StringTok}[1]{\textcolor[rgb]{0.25,0.44,0.63}{#1}}
\newcommand{\VariableTok}[1]{\textcolor[rgb]{0.10,0.09,0.49}{#1}}
\newcommand{\VerbatimStringTok}[1]{\textcolor[rgb]{0.25,0.44,0.63}{#1}}
\newcommand{\WarningTok}[1]{\textcolor[rgb]{0.38,0.63,0.69}{\textbf{\textit{#1}}}}
\usepackage{longtable,booktabs,array}
\usepackage{calc} % for calculating minipage widths
% Correct order of tables after \paragraph or \subparagraph
\usepackage{etoolbox}
\makeatletter
\patchcmd\longtable{\par}{\if@noskipsec\mbox{}\fi\par}{}{}
\makeatother
% Allow footnotes in longtable head/foot
\IfFileExists{footnotehyper.sty}{\usepackage{footnotehyper}}{\usepackage{footnote}}
\makesavenoteenv{longtable}
\setlength{\emergencystretch}{3em} % prevent overfull lines
\providecommand{\tightlist}{%
  \setlength{\itemsep}{0pt}\setlength{\parskip}{0pt}}
\setcounter{secnumdepth}{5}
\ifLuaTeX
\usepackage[bidi=basic]{babel}
\else
\usepackage[bidi=default]{babel}
\fi
\babelprovide[main,import]{chinese}
% get rid of language-specific shorthands (see #6817):
\let\LanguageShortHands\languageshorthands
\def\languageshorthands#1{}
\usepackage{booktabs,longtable,array,enumitem,xcolor,fancyhdr,lastpage}
\definecolor{BelBlue}{HTML}{1F4E79}
\pagestyle{fancy}\fancyhf{}\fancyhead[L]{\small 2003 BelPhO 10-2}\fancyhead[R]{\small 小行星球形空腔反演}\fancyfoot[C]{\small\thepage/\pageref{LastPage}}
\ifLuaTeX
  \usepackage{selnolig}  % disable illegal ligatures
\fi
\IfFileExists{bookmark.sty}{\usepackage{bookmark}}{\usepackage{hyperref}}
\IfFileExists{xurl.sty}{\usepackage{xurl}}{} % add URL line breaks if available
\urlstyle{same}
\hypersetup{
  pdftitle={2003 年白俄罗斯共和国物理奥林匹克决赛·10 年级第 2 题：小行星球形空腔反演},
  pdfauthor={中文竞赛精讲与独立复算版},
  pdflang={zh-CN},
  hidelinks,
  pdfcreator={LaTeX via pandoc}}

\title{2003 年白俄罗斯共和国物理奥林匹克决赛·10 年级第 2
题：小行星球形空腔反演}
\usepackage{etoolbox}
\makeatletter
\providecommand{\subtitle}[1]{% add subtitle to \maketitle
  \apptocmd{\@title}{\par {\large #1 \par}}{}{}
}
\makeatother
\subtitle{题源核验、完整建模、零跳步推导、官方答案复算与同类题迁移}
\author{中文竞赛精讲与独立复算版}
\date{2026-10-17}

\begin{document}
\maketitle

{
\setcounter{tocdepth}{3}
\tableofcontents
}
\hypertarget{ux4f7fux7528ux8bf4ux660eux4e0eux5b8cux6210ux58f0ux660e}{%
\section{使用说明与完成声明}\label{ux4f7fux7528ux8bf4ux660eux4e0eux5b8cux6210ux58f0ux660e}}

本文只围绕用户指定的数据库条目展开。数据库行把 2003 年中组理论题记为共 8
题；BelPhO 官方年度页依次列出 9 年级理论题 1---4 和 10 年级理论题
1---4，因此顺序号 6 对应官方 \textbf{10 年级第 2
题}。官方独立题页把它标为``2003 年共和国物理奥林匹克，10 年级，第 2
题''，站内难度 4/5；年度合订题卷第 3 页给出题面，年度解答第 6
页给出解法。

\textbf{建议文件名：}
\texttt{2003\_白俄罗斯共和国物理奥林匹克决赛\_10年级第2题\_小行星球形空腔反演\_完整精讲.md}

\textbf{来源定位：}

\begin{itemize}
\tightlist
\item
  官方年度页：\texttt{https://belpho.org/problems/final/2003/}
\item
  官方独立题页：\texttt{https://belpho.org/problems/final/2003/10-2/}
\item
  官方题卷：\texttt{https://belpho.org/problems/final/2003problems.pdf}，PDF
  第 3 页。
\item
  官方解答：\texttt{https://belpho.org/problems/final/2003solutions.pdf}，PDF
  第 6 页。
\item
  本文对官方公式和数值进行了独立复算；凡属本文补充的推导、误差分析和推广，均明确标为``独立推导''或``补充背景''。
\end{itemize}

\hypertarget{ux9898ux6e90ux96beux5ea6ux5206ux7c7bux4e0eux8badux7ec3ux76eeux6807}{%
\section{1.
题源、难度、分类与训练目标}\label{ux9898ux6e90ux96beux5ea6ux5206ux7c7bux4e0eux8badux7ec3ux76eeux6807}}

\begin{longtable}[]{@{}
  >{\raggedright\arraybackslash}p{(\columnwidth - 2\tabcolsep) * \real{0.5000}}
  >{\raggedright\arraybackslash}p{(\columnwidth - 2\tabcolsep) * \real{0.5000}}@{}}
\toprule\noalign{}
\begin{minipage}[b]{\linewidth}\raggedright
项目
\end{minipage} & \begin{minipage}[b]{\linewidth}\raggedright
内容
\end{minipage} \\
\midrule\noalign{}
\endhead
\bottomrule\noalign{}
\endlastfoot
赛事 & 白俄罗斯共和国物理奥林匹克全国决赛 \\
年份与组别 & 2003 年，10 年级 \\
官方题号 & 理论第 2 题，记作 10-2 \\
数据库顺序号 & 中组理论题第 6/8 题 \\
官方站难度 & 4/5 \\
本文统一难度 & 难，6/10 \\
分类标签 &
力学---万有引力；球壳定理；叠加原理；反问题；连续介质；无量纲化 \\
主考点 &
把空腔等效为``缺失质量球''；由两个对径测点反演空腔中心位置和半径；检查球外场公式的适用条件 \\
\end{longtable}

\hypertarget{ux4e3aux4ec0ux4e48ux662fux96be610}{%
\subsection{1.1
为什么是``难，6/10''}\label{ux4e3aux4ec0ux4e48ux662fux96be610}}

代数本身不长，真正的门槛有四个：

\begin{enumerate}
\def\labelenumi{\arabic{enumi}.}
\tightlist
\item
  看到``挖出球形空腔''时，必须想到``完整均匀球减去一个同密度小球''，而不是直接积分剩余不规则天体。
\item
  必须判断重力最小点、最大点与空腔中心共线，并正确写出到空腔中心的两个距离
  \(a\) 与 \(2R-a\)。
\item
  必须区分缺失质量按 \(r^3\) 缩放和场强按距离平方衰减。
\item
  需要从两个异常值中消去未知密度、万有引力常量和空腔半径，先反演位置，再反演大小。
\end{enumerate}

它属于国家决赛常规难题，但推导链清晰，不需要国际奥赛极难题级别的多次创造性突破，因此定为
6/10，而不是 8---10 分。

\hypertarget{ux5b98ux65b9ux539fux9898ux7684ux5408ux89c4ux5448ux73b0ux4e0eux5b8cux6574ux4e2dux6587ux9898ux610f}{%
\section{2.
官方原题的合规呈现与完整中文题意}\label{ux5b98ux65b9ux539fux9898ux7684ux5408ux89c4ux5448ux73b0ux4e0eux5b8cux6574ux4e2dux6587ux9898ux610f}}

受公开材料的转载限制，本文不逐字大段复制俄文题面，只保留一个短节选并给出不省略条件的中文重述。完整俄文原文请直接查看官方
PDF 第 3 页。

\textbf{俄文短节选：} ``В глубинах Вселенной был обнаружен однородный
астероид сферической формы радиуса R \ldots{}''

\hypertarget{ux5b8cux6574ux4e2dux6587ux9898ux610fux91cdux8ff0}{%
\subsection{2.1
完整中文题意重述}\label{ux5b8cux6574ux4e2dux6587ux9898ux610fux91cdux8ff0}}

在宇宙深处发现了一颗密度均匀、半径为 \(R\)
的球形小行星。最初用高精度重力仪测量时，表面各点的重力加速度大小都相同，记为
\(g_0\)。

后来在小行星内部采矿，形成一个没有与表面相交的球形空腔。再次测量发现：

\begin{itemize}
\tightlist
\item
  表面某点 \(A\) 的重力加速度最小，\(g_{\min}=0.938g_0\)，即相对原值减少
  \(η_1=6.2\%\)；
\item
  与 \(A\) 直径相对的表面点
  \(B\)，重力加速度最大，\(g_{\max}=0.993g_0\)，即相对原值减少
  \(η_2=0.70\%\)。
\end{itemize}

根据这些测量数据，确定球形空腔中心的位置、中心在表面下的埋深
\(a\)，以及空腔半径 \(r\)。

\hypertarget{ux56feux5f62ux7b26ux53f7ux4e0eux65b9ux5411ux7ea6ux5b9a}{%
\section{3.
图形、符号与方向约定}\label{ux56feux5f62ux7b26ux53f7ux4e0eux65b9ux5411ux7ea6ux5b9a}}

沿直径 \(AB\) 建立一维几何关系：

\begin{Shaded}
\begin{Highlighting}[]
\NormalTok{A {-}{-}{-}{-}{-}{-}{-}{-} C {-}{-}{-}{-}{-}{-}{-}{-} O {-}{-}{-}{-}{-}{-}{-}{-}{-}{-}{-}{-}{-}{-}{-}{-} B}
\NormalTok{|\textless{}{-}{-} a {-}{-}\textgreater{}|\textless{}{-}{-} R{-}a {-}{-}\textgreater{}|\textless{}{-}{-}{-}{-}{-}{-} R {-}{-}{-}{-}{-}{-}\textgreater{}|}


\NormalTok{A：重力最小点}
\NormalTok{B：重力最大点，且与 A 对径}
\NormalTok{O：原小行星球心}
\NormalTok{C：空腔中心}
\NormalTok{AC=a，CB=2R{-}a，OC=R{-}a}
\NormalTok{空腔半径为 r}
\end{Highlighting}
\end{Shaded}

符号统一如下：

\begin{itemize}
\tightlist
\item
  \(R\)：原小行星半径。
\item
  \(ρ\)：原小行星物质的均匀密度。
\item
  \(M=(4π/3)ρR^3\)：原完整小行星质量。
\item
  \(g_0\)：未采矿时表面重力大小。
\item
  \(r\)：空腔半径。
\item
  \(m=(4π/3)ρr^3\)：被挖去物质的质量。
\item
  \(a=AC\)：空腔中心到最近表面点 \(A\) 的距离，也就是题目所称中心埋深。
\item
  \(η_1=0.062\)，\(η_2=0.007\)：两处相对重力减少量。
\end{itemize}

【卡点提醒】\(a\) 是``表面点 \(A\)
到空腔中心''的距离，不是空腔上边界到表面的净覆岩厚度。后者应为 \(a-r\)。

\hypertarget{ux8bfbux9898ux62c6ux89e3}{%
\section{4. 读题拆解}\label{ux8bfbux9898ux62c6ux89e3}}

\hypertarget{ux5927ux767dux8bddux63cfux8ff0}{%
\subsection{4.1 大白话描述}\label{ux5927ux767dux8bddux63cfux8ff0}}

原来是一个完全球对称的小行星，所以表面处处重力相同。挖掉一团球形物质后，离空腔最近的表面点少受的引力最多，离空腔最远的对径点少受的引力最少。题目给了这两个``缺失引力''的大小，要我们反推出空腔在哪儿、有多大。

\hypertarget{ux663eux5f0fux6761ux4ef6}{%
\subsection{4.2 显式条件}\label{ux663eux5f0fux6761ux4ef6}}

\begin{itemize}
\tightlist
\item
  原天体是均匀实心球，半径 \(R\)。
\item
  空腔为球形，且完全位于小行星内部。
\item
  \(A\)、\(B\) 为对径点。
\item
  \(g_{\min}=0.938g_0\)，\(g_{\max}=0.993g_0\)。
\item
  求空腔中心位置、埋深 \(a\) 和半径 \(r\)。
\end{itemize}

\hypertarget{ux9690ux85cfux6761ux4ef6ux4e0eux547dux9898ux7ed3ux6784}{%
\subsection{4.3
隐藏条件与命题结构}\label{ux9690ux85cfux6761ux4ef6ux4e0eux547dux9898ux7ed3ux6784}}

\begin{enumerate}
\def\labelenumi{\arabic{enumi}.}
\tightlist
\item
  \textbf{重力最小点必在空腔最近侧。}
  缺失质量造成的场修正最大处，就是到空腔中心距离最小的表面点。
\item
  \textbf{两个极值点与空腔中心共线。}
  对球形缺失质量，外部场只依赖到球心的距离；距离在表面球上取得最小和最大值的位置，正是通过两个球心的直线与表面的两个交点。
\item
  \textbf{\(A\)、\(C\)、\(O\)、\(B\) 的顺序是 \(A-C-O-B\)。} 因为 \(A\)
  是最近点，空腔中心位于 \(A\) 一侧。
\item
  \textbf{\(A\)、\(B\) 都在``被挖去的小球''的外部。} 空腔不出露，故
  \(a>r\)；远点显然也在外部。因此缺失球在两点的场可按点质量处理。
\item
  \textbf{\(g_{\max}\) 仍小于 \(g_0\)。}
  ``最大''只是在采矿后的各表面点之间比较，不代表它超过采矿前的正常值。
\end{enumerate}

\hypertarget{ux547dux9898ux4ebaux771fux6b63ux60f3ux8003ux4ec0ux4e48}{%
\subsection{4.4
命题人真正想考什么}\label{ux547dux9898ux4ebaux771fux6b63ux60f3ux8003ux4ec0ux4e48}}

不是考繁琐积分，而是考是否会把复杂几何改写成线性叠加：

\[\text{有空腔的小行星}=\text{完整均匀球}-\text{空腔位置处的均匀小球}.\]

这就是``负密度球''或``补回缺失质量''的方法。负密度只是计算记账，不表示真实存在负质量。

\hypertarget{ux7269ux7406ux5efaux6a21ux516dux6b65ux6cd5}{%
\section{5.
物理建模六步法}\label{ux7269ux7406ux5efaux6a21ux516dux6b65ux6cd5}}

\hypertarget{ux7b2cux4e00ux6b65ux7814ux7a76ux5bf9ux8c61ux4e0eux7cfbux7edfux8fb9ux754c}{%
\subsection{第一步：研究对象与系统边界}\label{ux7b2cux4e00ux6b65ux7814ux7a76ux5bf9ux8c61ux4e0eux7cfbux7edfux8fb9ux754c}}

研究表面点 \(A\)、\(B\)
处单位质量试探物体受到的引力场。引力源拆成：完整大球与一个位于
\(C\)、质量为 \(m\) 的``被移除小球''。

\hypertarget{ux7b2cux4e8cux6b65ux53c2ux8003ux7cfbux548cux5750ux6807}{%
\subsection{第二步：参考系和坐标}\label{ux7b2cux4e8cux6b65ux53c2ux8003ux7cfbux548cux5750ux6807}}

取小行星质心附近的非旋转参考系，沿 \(A\to B\)
建立坐标轴。由于本题只比较静态引力，不涉及惯性力或运动学过程。

\hypertarget{ux7b2cux4e09ux6b65ux72b6ux6001ux4e0eux8fc7ux7a0b}{%
\subsection{第三步：状态与过程}\label{ux7b2cux4e09ux6b65ux72b6ux6001ux4e0eux8fc7ux7a0b}}

这是采矿前后两个静态引力场的比较：

\begin{itemize}
\tightlist
\item
  采矿前：完整均匀球，表面场大小为 \(g_0\)。
\item
  采矿后：完整球场减去缺失球原本产生的场。
\end{itemize}

\hypertarget{ux7b2cux56dbux6b65ux76f8ux4e92ux4f5cux7528ux4e0eux51e0ux4f55ux7ea6ux675f}{%
\subsection{第四步：相互作用与几何约束}\label{ux7b2cux56dbux6b65ux76f8ux4e92ux4f5cux7528ux4e0eux51e0ux4f55ux7ea6ux675f}}

只有牛顿引力。几何约束为
\(AC=a\)、\(CB=2R-a\)、\(OC=R-a\)，并且空腔完全在内部，最终必须检验
\(a-r>0\)。

\hypertarget{ux7b2cux4e94ux6b65ux9002ux7528ux5b9aux5f8bux4e0eux6210ux7acbux6761ux4ef6}{%
\subsection{第五步：适用定律与成立条件}\label{ux7b2cux4e94ux6b65ux9002ux7528ux5b9aux5f8bux4e0eux6210ux7acbux6761ux4ef6}}

\begin{itemize}
\tightlist
\item
  牛顿万有引力定律。
\item
  引力场线性叠加原理。
\item
  均匀球的球壳定理：球外场等效于质量集中在球心。
\item
  均匀实心球表面场 \(g_0=GM/R^2\)。
\end{itemize}

\hypertarget{ux7b2cux516dux6b65ux7ed3ux679cux81eaux68c0ux8ba1ux5212}{%
\subsection{第六步：结果自检计划}\label{ux7b2cux516dux6b65ux7ed3ux679cux81eaux68c0ux8ba1ux5212}}

至少完成：

\begin{enumerate}
\def\labelenumi{\arabic{enumi}.}
\tightlist
\item
  量纲检查；
\item
  反代两处异常比例；
\item
  验证 \(a>r\)，空腔没有穿出表面；
\item
  验证 \(η_1>η_2\) 对应空腔中心偏向 \(A\)；
\item
  检查极限 \(η_1=η_2\) 是否给出同心空腔。
\end{enumerate}

\hypertarget{ux7a81ux7834ux53e3ux5206ux6790}{%
\section{6. 突破口分析}\label{ux7a81ux7834ux53e3ux5206ux6790}}

\hypertarget{ux7b2cux4e00ux53cdux5e94ux5148ux7b97ux5c11ux6389ux591aux5c11ux5f15ux529bux4e0dux8981ux76f4ux63a5ux7b97ux5269ux4f59ux5929ux4f53}{%
\subsection{6.1
第一反应：先算``少掉多少引力''，不要直接算剩余天体}\label{ux7b2cux4e00ux53cdux5e94ux5148ux7b97ux5c11ux6389ux591aux5c11ux5f15ux529bux4e0dux8981ux76f4ux63a5ux7b97ux5269ux4f59ux5929ux4f53}}

有空腔后，整体不再球对称。若直接对剩余体积积分，会把一道两式消元题做成三维积分。更稳定的路线是把被挖去的物质``补回去''：

\[g_0=g_{\text{after}}+g_{\text{missing sphere}}.\]

于是两处测量给出的正好是缺失球原本会产生的场。

\hypertarget{ux4e3aux4ec0ux4e48ux5148ux53d6ux6bd4ux503c}{%
\subsection{6.2
为什么先取比值}\label{ux4e3aux4ec0ux4e48ux5148ux53d6ux6bd4ux503c}}

在 \(A\)、\(B\) 两点，缺失球的质量 \(m\) 相同，万有引力常量 \(G\)
相同。两式相除会同时消去 \(G\)、\(m\)、\(ρ\) 和
\(r\)，只留下几何距离比，所以应当先求位置 \(a\)。

\hypertarget{ux8003ux573aux51e0ux5341ux79d2ux8bc6ux522bux903bux8f91}{%
\subsection{6.3
考场几十秒识别逻辑}\label{ux8003ux573aux51e0ux5341ux79d2ux8bc6ux522bux903bux8f91}}

看到以下组合时立即想到``负质量球 + 比值消元''：

\begin{itemize}
\tightlist
\item
  均匀球；
\item
  内部球形空腔；
\item
  外部或表面场测量；
\item
  两个不同位置的异常值；
\item
  求空腔位置和大小。
\end{itemize}

\hypertarget{ux770bux4f3cux5408ux7406ux4f46ux8d70ux4e0dux901aux7684ux8defux7ebf}{%
\subsection{6.4
看似合理但走不通的路线}\label{ux770bux4f3cux5408ux7406ux4f46ux8d70ux4e0dux901aux7684ux8defux7ebf}}

\begin{itemize}
\tightlist
\item
  \textbf{把有洞天体整体当作地心点质量。}
  错在剩余质量分布不再球对称；只知道总质量不足以确定表面不同位置的场。
\item
  \textbf{把距离写成 \(a-r\)。} 平方反比律中的源点是空腔中心
  \(C\)，不是空腔上边界。
\item
  \textbf{直接令 \(a=R/2\)、\(r=R/4\)。}
  这只是数值近似，不能在推导前当作精确几何条件。
\end{itemize}

\hypertarget{ux7b2cux4e00ux79cdux5b8cux6574ux89e3ux6cd5ux8865ux56deux7f3aux5931ux8d28ux91cf}{%
\section{7.
第一种完整解法：补回缺失质量}\label{ux7b2cux4e00ux79cdux5b8cux6574ux89e3ux6cd5ux8865ux56deux7f3aux5931ux8d28ux91cf}}

\hypertarget{ux539fux5b8cux6574ux5c0fux884cux661fux7684ux8868ux9762ux91cdux529b}{%
\subsection{7.1
原完整小行星的表面重力}\label{ux539fux5b8cux6574ux5c0fux884cux661fux7684ux8868ux9762ux91cdux529b}}

原质量为

\[M=(4π/3)ρR^3.\]

表面重力为

\[g_0=GM/R^2=(4πGρ/3)R.\]

【思路由来】把 \(g_0\) 写成含 \(ρ\)
的形式，是为了稍后与缺失球的场相除，让 \(G\) 与 \(ρ\) 自动消去。

\hypertarget{ux88abux6316ux53bbux7269ux8d28ux7684ux8d28ux91cf}{%
\subsection{7.2
被挖去物质的质量}\label{ux88abux6316ux53bbux7269ux8d28ux7684ux8d28ux91cf}}

球形空腔半径为 \(r\)，原本该区域装有同密度物质，所以缺失质量为

\[m=(4π/3)ρr^3.\]

【卡点提醒】质量按半径三次方缩放；若误写成
\(r^2\)，后面即使代数无误，物理量纲也会失败。

\hypertarget{ux70b9-a-ux7684ux91cdux529bux5f02ux5e38}{%
\subsection{7.3 点 A
的重力异常}\label{ux70b9-a-ux7684ux91cdux529bux5f02ux5e38}}

采矿后点 \(A\) 的重力减少量为

\[Δg_A=g_0-g_{\min}=(1-0.938)g_0=0.062g_0=η_1g_0.\]

点 \(A\) 到缺失球中心的距离为 \(a\)。因 \(A\) 在缺失球外部，球壳定理给出

\[Δg_A=Gm/a^2.\]

因此

\[η_1g_0=Gm/a^2. \tag{1}\]

\hypertarget{ux70b9-b-ux7684ux91cdux529bux5f02ux5e38}{%
\subsection{7.4 点 B
的重力异常}\label{ux70b9-b-ux7684ux91cdux529bux5f02ux5e38}}

点 \(B\) 的减少量为

\[Δg_B=g_0-g_{\max}=(1-0.993)g_0=0.007g_0=η_2g_0.\]

从 \(B\) 到空腔中心的距离是

\[BC=BA-AC=2R-a.\]

所以

\[η_2g_0=Gm/(2R-a)^2. \tag{2}\]

\hypertarget{ux4e24ux5f0fux76f8ux9664ux5148ux6c42ux4e2dux5fc3ux57cbux6df1}{%
\subsection{7.5
两式相除，先求中心埋深}\label{ux4e24ux5f0fux76f8ux9664ux5148ux6c42ux4e2dux5fc3ux57cbux6df1}}

用式 (1) 除以式 (2)：

\[η_1/η_2=(2R-a)^2/a^2.\]

所有量均为正，取正平方根：

\[√(η_1/η_2)=(2R-a)/a.\]

逐步移项：

\[a√η_1=(2R-a)√η_2,\]

\[a√η_1+a√η_2=2R√η_2,\]

\[a(√η_1+√η_2)=2R√η_2.\]

于是

\[a/R=2√η_2/(√η_1+√η_2). \tag{3}\]

代入 \(η_1=0.062\)、\(η_2=0.007\)：

\[√η_1≈0.248998,\]

\[√η_2≈0.0836660,\]

\[a/R≈2(0.0836660)/(0.248998+0.0836660)=0.503006.\]

所以

\[a≈0.503R.\]

空腔中心相对小行星球心的位置为

\[OC=R-a≈0.496994R,\]

方向从 \(O\) 指向 \(A\)。

\hypertarget{ux518dux6c42ux7a7aux8154ux534aux5f84}{%
\subsection{7.6
再求空腔半径}\label{ux518dux6c42ux7a7aux8154ux534aux5f84}}

由式 (1)：

\[η_1g_0=Gm/a^2.\]

代入 \(g_0=(4πGρ/3)R\) 和 \(m=(4πρ/3)r^3\)：

\[η_1(4πGρ/3)R=G(4πρ/3)r^3/a^2.\]

两边约去共同因子 \(4πGρ/3\)：

\[η_1R=r^3/a^2.\]

因此

\[r^3=η_1Ra^2.\]

除以 \(R^3\)：

\[(r/R)^3=η_1(a/R)^2.\]

代入式 (3)，可得闭式

\[r/R=[4η_1η_2/(√η_1+√η_2)^2]^{1/3}. \tag{4}\]

数值为

\[(r/R)^3≈0.01568694,\]

\[r/R≈0.2503299.\]

所以

\[r≈0.250R.\]

\hypertarget{ux7b2cux4e00ux89e3ux6cd5ux6700ux7ec8ux7ed3ux8bba}{%
\subsection{7.7
第一解法最终结论}\label{ux7b2cux4e00ux89e3ux6cd5ux6700ux7ec8ux7ed3ux8bba}}

\begin{itemize}
\tightlist
\item
  空腔中心位于直径 \(AO\) 上、靠近 \(A\) 的一侧。
\item
  中心埋深：\(a=0.503006R≈0.503R≈R/2\)。
\item
  相对球心偏移：\(OC=0.496994R\)，方向指向 \(A\)。
\item
  空腔半径：\(r=0.250330R≈0.250R≈R/4\)。
\item
  空腔最靠近表面的边界深度：\(a-r≈0.252676R\)。
\end{itemize}

【卡点提醒】\(R/2\) 与 \(R/4\) 只是三位有效数字下的近似，不是精确答案。

\hypertarget{ux7b2cux4e8cux79cdux5b8cux6574ux89e3ux6cd5ux65e0ux91cfux7eb2ux53cdux6f14}{%
\section{8.
第二种完整解法：无量纲反演}\label{ux7b2cux4e8cux79cdux5b8cux6574ux89e3ux6cd5ux65e0ux91cfux7eb2ux53cdux6f14}}

这条路线从一开始就把未知量压缩成两个无量纲数，适合考场快速复算。

定义

\[x=a/R,\]

\[f=m/M=(r/R)^3.\]

完整球表面场是 \(g_0=GM/R^2\)。缺失球在 \(A\) 处造成的相对异常为

\[η_1=(Gm/a^2)/(GM/R^2)=f/x^2. \tag{5}\]

在 \(B\) 处：

\[η_2=(Gm/(2R-a)^2)/(GM/R^2)=f/(2-x)^2. \tag{6}\]

式 (5)、(6) 开平方：

\[√f=x√η_1,\]

\[√f=(2-x)√η_2.\]

令右边相等：

\[x√η_1=(2-x)√η_2,\]

从而

\[x=2√η_2/(√η_1+√η_2)=0.503006.\]

再由 \(f=η_1x^2\)：

\[f=0.062(0.503006)^2=0.01568694.\]

因为 \(f=(r/R)^3\)，故

\[r/R=f^{1/3}=0.2503299.\]

\hypertarget{ux89e3ux6cd5ux5bf9ux6bd4}{%
\subsection{8.1 解法对比}\label{ux89e3ux6cd5ux5bf9ux6bd4}}

\begin{itemize}
\tightlist
\item
  第一解法最适合首次学习：清楚展示质量、密度、球壳定理和``补回物质''的物理含义。
\item
  第二解法最适合考场：未知量只有
  \(x\)、\(f\)，方程直接揭示``位置由异常比决定，大小由异常绝对量决定''。
\item
  两者组织方式不同，但物理根基相同；不能把换符号后的同一推导冒充完全独立的第三种理论。
\end{itemize}

\hypertarget{ux5b98ux65b9ux89e3ux7b54ux4e0eux72ecux7acbux590dux7b97ux5bf9ux7167}{%
\section{9.
官方解答与独立复算对照}\label{ux5b98ux65b9ux89e3ux7b54ux4e0eux72ecux7acbux590dux7b97ux5bf9ux7167}}

官方解答采用``把挖去的物质补回去''的路线，先由两处异常之比求 \(a\)，再求
\(r\)。官方给出的最终数值为

\[a=0.503R≈R/2,\]

\[r=0.250R≈R/4.\]

独立复算保留更多位数：

\[a=0.5030060602R,\]

\[r=0.2503298898R.\]

\begin{longtable}[]{@{}
  >{\raggedright\arraybackslash}p{(\columnwidth - 8\tabcolsep) * \real{0.1667}}
  >{\raggedleft\arraybackslash}p{(\columnwidth - 8\tabcolsep) * \real{0.2222}}
  >{\raggedleft\arraybackslash}p{(\columnwidth - 8\tabcolsep) * \real{0.2222}}
  >{\raggedleft\arraybackslash}p{(\columnwidth - 8\tabcolsep) * \real{0.2222}}
  >{\raggedright\arraybackslash}p{(\columnwidth - 8\tabcolsep) * \real{0.1667}}@{}}
\toprule\noalign{}
\begin{minipage}[b]{\linewidth}\raggedright
项目
\end{minipage} & \begin{minipage}[b]{\linewidth}\raggedleft
官方值
\end{minipage} & \begin{minipage}[b]{\linewidth}\raggedleft
独立复算值
\end{minipage} & \begin{minipage}[b]{\linewidth}\raggedleft
差异
\end{minipage} & \begin{minipage}[b]{\linewidth}\raggedright
结论
\end{minipage} \\
\midrule\noalign{}
\endhead
\bottomrule\noalign{}
\endlastfoot
中心埋深 \(a/R\) & 0.503 & 0.5030060602 & 约 \(6.1×10^{-6}\) &
正常舍入，一致 \\
空腔半径 \(r/R\) & 0.250 & 0.2503298898 & 约 \(3.30×10^{-4}\) &
官方保留三位小数，一致 \\
近似口径 & \(R/2\)、\(R/4\) & 仅为近似 & --- & 不应写成精确等号 \\
\end{longtable}

\hypertarget{ux5b98ux65b9ux89e3ux7b54ux4e2dux7684ux4e00ux4e2aux91cfux7eb2ux7b14ux8bef}{%
\subsection{9.1
官方解答中的一个量纲笔误}\label{ux5b98ux65b9ux89e3ux7b54ux4e2dux7684ux4e00ux4e2aux91cfux7eb2ux7b14ux8bef}}

官方解答在写点 \(A\) 的异常场时包含密度因子 \(ρ\)；紧接着单独展示点
\(B\) 异常场的那一行，排版文本漏掉了
\(ρ\)，写成的量纲不再是加速度。但随后列出的联立方程又恢复了
\(ρ\)，最后公式与数值也都按正确表达式计算。

正确表达式应为

\[Δg_B=G(4πρr^3/3)/(2R-a)^2.\]

因此这是局部排版笔误，不影响官方最终答案。本文同时保留``官方原貌存在漏因子''和``严格复算采用正确式子''，不作静默修改。

\hypertarget{ux56dbux91cdux6b63ux786eux6027ux68c0ux67e5}{%
\section{10.
四重正确性检查}\label{ux56dbux91cdux6b63ux786eux6027ux68c0ux67e5}}

\hypertarget{ux6587ux672cux68c0ux67e5}{%
\subsection{10.1 文本检查}\label{ux6587ux672cux68c0ux67e5}}

\begin{itemize}
\tightlist
\item
  题目给的是 \(g_{\min}=0.938g_0\) 与
  \(g_{\max}=0.993g_0\)，相对减少量分别为 0.062 与 0.007。
\item
  \(B\) 与 \(A\) 对径，所以 \(AB=2R\)。
\item
  所求 \(a\) 是中心埋深 \(AC\)，而不是空腔边界深度。
\end{itemize}

\hypertarget{ux6570ux5b66ux53cdux4ee3}{%
\subsection{10.2 数学反代}\label{ux6570ux5b66ux53cdux4ee3}}

取

\[x=a/R=0.5030060602,\]

\[f=(r/R)^3=0.01568693599.\]

则

\[f/x^2=0.0620000000=η_1,\]

\[f/(2-x)^2=0.0070000000=η_2.\]

两处测量均被精确复现到显示精度。

\hypertarget{ux7269ux7406ux68c0ux67e5}{%
\subsection{10.3 物理检查}\label{ux7269ux7406ux68c0ux67e5}}

\hypertarget{ux91cfux7eb2}{%
\subsubsection{量纲}\label{ux91cfux7eb2}}

\(Gm/a^2\) 的量纲是加速度；\(η_1\)、\(η_2\)、\(a/R\)、\(r/R\) 均无量纲。

\hypertarget{ux65b9ux5411ux548cux5927ux5c0f}{%
\subsubsection{方向和大小}\label{ux65b9ux5411ux548cux5927ux5c0f}}

\(η_1>η_2\)，说明 \(A\) 离空腔中心更近，故空腔偏向 \(A\)
一侧。\(g_{\max}=0.993g_0\) 虽是采矿后表面最大值，仍小于原来的
\(g_0\)，符合``移走正质量只会减弱引力''的判断。

\hypertarget{ux7a7aux8154ux786eux5b9eux5b8cux5168ux4f4dux4e8eux5185ux90e8}{%
\subsubsection{空腔确实完全位于内部}\label{ux7a7aux8154ux786eux5b9eux5b8cux5168ux4f4dux4e8eux5185ux90e8}}

最近侧净覆岩厚度为

\[a-r=(0.503006-0.250330)R=0.252676R>0.\]

远侧到表面的剩余距离为

\[2R-a-r=1.246664R>0.\]

所以所得几何与``空腔不出露''自洽。

\hypertarget{ux7f3aux5931ux8d28ux91cfux6bd4ux4f8b}{%
\subsubsection{缺失质量比例}\label{ux7f3aux5931ux8d28ux91cfux6bd4ux4f8b}}

\[m/M=(r/R)^3≈0.0156869.\]

只移除了约 \(1.57\%\) 的总质量，却使最近表面点的重力减少
\(6.2\%\)。这不矛盾，因为缺失质量离 \(A\) 的距离仅约
\(0.503R\)，远小于完整天体质量中心到表面的距离 \(R\)。

\hypertarget{ux6781ux9650ux68c0ux67e5}{%
\subsection{10.4 极限检查}\label{ux6781ux9650ux68c0ux67e5}}

若 \(η_1=η_2\)，式 (3) 给出

\[a/R=1,\]

即空腔中心与大球中心重合。此时所有表面点的重力减少量相同，不会出现方向性差异，完全符合球对称性。

若 \(η_2\to0\) 而 \(η_1\) 保持有限，则 \(a/R\to0\)，意味着空腔中心逼近点
\(A\)；远点异常趋于很小，也符合直觉。不过空腔仍需满足
\(r<a\)，所以实际参数不能无限逼近该极限。

\hypertarget{ux8868ux8fbeux68c0ux67e5}{%
\subsection{10.5 表达检查}\label{ux8868ux8fbeux68c0ux67e5}}

\begin{itemize}
\tightlist
\item
  所有最终比值都注明相对于 \(R\)。
\item
  近似号与等号严格区分。
\item
  官方值、独立复算值、补充推广分别标注。
\end{itemize}

\hypertarget{ux5e38ux89c1ux8befux533a}{%
\section{11. 常见误区}\label{ux5e38ux89c1ux8befux533a}}

\hypertarget{ux8befux533aux4e00ux7528-a-r-ux4ee3ux66ff-a}{%
\subsection{\texorpdfstring{误区一：用 \(a-r\) 代替
\(a\)}{误区一：用 a-r 代替 a}}\label{ux8befux533aux4e00ux7528-a-r-ux4ee3ux66ff-a}}

\begin{itemize}
\tightlist
\item
  \textbf{错误表现：} 写 \(Δg_A=Gm/(a-r)^2\)。
\item
  \textbf{错误原因：}
  把``到空腔上边界的距离''误当成``到等效点质量的距离''。
\item
  \textbf{如何自检：}
  球壳定理把球外场等效为质量集中在球心，不会集中在最近表面点。
\item
  \textbf{正确改法：} 源点放在空腔中心 \(C\)，距离为 \(AC=a\)。
\end{itemize}

\hypertarget{ux8befux533aux4e8cux628aux8d28ux91cfux6bd4ux4f8bux5199ux6210ux9762ux79efux6bd4ux4f8b}{%
\subsection{误区二：把质量比例写成面积比例}\label{ux8befux533aux4e8cux628aux8d28ux91cfux6bd4ux4f8bux5199ux6210ux9762ux79efux6bd4ux4f8b}}

\begin{itemize}
\tightlist
\item
  \textbf{错误表现：} 令 \(m/M=(r/R)^2\)。
\item
  \textbf{错误原因：} 忘记三维球体的体积按半径三次方缩放。
\item
  \textbf{如何自检：} 密度乘面积不是质量，量纲缺少一个长度。
\item
  \textbf{正确改法：} \(m/M=(r/R)^3\)。
\end{itemize}

\hypertarget{ux8befux533aux4e09ux8ba4ux4e3a-g_max-ux5927ux4e8e-g_0}{%
\subsection{\texorpdfstring{误区三：认为 \(g_{\max}\) 大于
\(g_0\)}{误区三：认为 g\_\{\textbackslash max\} 大于 g\_0}}\label{ux8befux533aux4e09ux8ba4ux4e3a-g_max-ux5927ux4e8e-g_0}}

\begin{itemize}
\tightlist
\item
  \textbf{错误表现：} 写 \(g_{\max}-g_0=0.007g_0\)。
\item
  \textbf{错误原因：}
  混淆``采矿后各点中的最大值''与``相对采矿前是否增大''。
\item
  \textbf{如何自检：}
  移走正质量后，不可能在同一点增加其朝向原天体的引力。
\item
  \textbf{正确改法：} \(g_0-g_{\max}=0.007g_0\)。
\end{itemize}

\hypertarget{ux8befux533aux56dbux628aux6709ux6d1eux5929ux4f53ux5f53ux4f5cux603bux8d28ux91cfux96c6ux4e2dux5728-o}{%
\subsection{误区四：把有洞天体当作总质量集中在
O}\label{ux8befux533aux56dbux628aux6709ux6d1eux5929ux4f53ux5f53ux4f5cux603bux8d28ux91cfux96c6ux4e2dux5728-o}}

\begin{itemize}
\tightlist
\item
  \textbf{错误表现：} 用剩余总质量 \(M-m\) 写所有表面点的场都等于
  \(G(M-m)/R^2\)。
\item
  \textbf{错误原因：} 球壳定理要求球对称；偏心空腔破坏了球对称。
\item
  \textbf{如何自检：}
  这种错误模型会预言所有表面点重力相同，与题目给出的两个极值直接冲突。
\item
  \textbf{正确改法：} 完整大球场与缺失小球场分别计算，再作向量叠加。
\end{itemize}

\hypertarget{ux8befux533aux4e94ux5728ux4efbux610fux8868ux9762ux70b9ux53eaux51cfux573aux5f3aux5927ux5c0f}{%
\subsection{误区五：在任意表面点只减场强大小}\label{ux8befux533aux4e94ux5728ux4efbux610fux8868ux9762ux70b9ux53eaux51cfux573aux5f3aux5927ux5c0f}}

\begin{itemize}
\tightlist
\item
  \textbf{错误表现：} 对不在轴线上的点写 \(g=g_0-Gm/s^2\)。
\item
  \textbf{错误原因：} 两个场一般不共线，必须按向量相加。
\item
  \textbf{如何自检：} 只有 \(A\)、\(B\)
  这两个轴上点，完整球场与缺失球场才共线。
\item
  \textbf{正确改法：} 一般位置使用分量或向量公式。
\end{itemize}

\hypertarget{ux8befux533aux516dux628aux8fd1ux4f3cux5206ux6570ux5f53ux7cbeux786eux51e0ux4f55}{%
\subsection{误区六：把近似分数当精确几何}\label{ux8befux533aux516dux628aux8fd1ux4f3cux5206ux6570ux5f53ux7cbeux786eux51e0ux4f55}}

\begin{itemize}
\tightlist
\item
  \textbf{错误表现：} 直接宣布 \(a=R/2\)、\(r=R/4\)。
\item
  \textbf{错误原因：} 只看三位小数的漂亮形式，没有保留原始测量数据。
\item
  \textbf{如何自检：} 取 \(a=R/2\)、\(r=R/4\)
  反代，得到的异常不会同时精确等于 0.062 与 0.007。
\item
  \textbf{正确改法：} 先写 \(0.503006R\) 与
  \(0.250330R\)，最后再说明近似于 \(R/2\)、\(R/4\)。
\end{itemize}

\hypertarget{ux4e00ux822cux5316ux4e0eux9ad8ux7b49ux7269ux7406ux80ccux666f}{%
\section{12.
一般化与高等物理背景}\label{ux4e00ux822cux5316ux4e0eux9ad8ux7b49ux7269ux7406ux80ccux666f}}

\hypertarget{ux5bc6ux5ea6ux5deeux4e0dux7b49ux4e8eux539fux5bc6ux5ea6ux65f6}{%
\subsection{12.1
密度差不等于原密度时}\label{ux5bc6ux5ea6ux5deeux4e0dux7b49ux4e8eux539fux5bc6ux5ea6ux65f6}}

若空腔不是全空，而是被密度 \(ρ_c\) 的物质填充，定义密度亏损比例

\[κ=(ρ-ρ_c)/ρ.\]

异常方程变为

\[η_1=κ(r/R)^3/(a/R)^2,\]

\[η_2=κ(r/R)^3/(2-a/R)^2.\]

两式相除后 \(κ\) 仍消去，所以中心位置 \(a\) 不变；但半径变为

\[r/R=[4η_1η_2/(κ(√η_1+√η_2)^2)]^{1/3}.\]

若 \(κ\) 未知，仅凭两点数据只能确定组合
\(κr^3\)，不能分别唯一确定密度差与半径。这是反问题中的``参数不可辨识性''。

\hypertarget{ux4e00ux822cux8868ux9762ux70b9ux7684ux5411ux91cfux573a}{%
\subsection{12.2
一般表面点的向量场}\label{ux4e00ux822cux8868ux9762ux70b9ux7684ux5411ux91cfux573a}}

设小行星中心为原点，空腔中心位置向量为 \(\boldsymbol c\)，表面点位置为
\(\boldsymbol x\)，\(|\boldsymbol x|=R\)。则采矿后场为

\[\boldsymbol g(\boldsymbol x)=-GM\boldsymbol x/R^3+Gm(\boldsymbol x-\boldsymbol c)/|\boldsymbol x-\boldsymbol c|^3.\]

第二项带正号，是因为我们从完整球场中减去了``被挖去正质量球''原本产生的、指向
\(C\) 的场。沿 \(OC\)
轴的近点与远点正是表面场异常的极值候选；离轴点必须保留向量方向。

\hypertarget{ux504fux5fc3ux7403ux5f62ux7a7aux8154ux5185ux90e8ux4e3aux4f55ux662fux5300ux5f3aux573a}{%
\subsection{12.3
偏心球形空腔内部为何是匀强场}\label{ux504fux5fc3ux7403ux5f62ux7a7aux8154ux5185ux90e8ux4e3aux4f55ux662fux5300ux5f3aux573a}}

对完整均匀大球，球内某点的位置向量为 \(\boldsymbol x\) 时

\[\boldsymbol g_1=-(4πGρ/3)\boldsymbol x.\]

空腔可视为叠加一个中心在 \(\boldsymbol c\)、密度为 \(-ρ\)
的小球。空腔内部相对小球中心的位置向量为
\(\boldsymbol x-\boldsymbol c\)，负密度球产生

\[\boldsymbol g_2=+(4πGρ/3)(\boldsymbol x-\boldsymbol c).\]

相加得到

\[\boldsymbol g=-(4πGρ/3)\boldsymbol c,\]

与空腔内具体位置 \(\boldsymbol x\)
无关。因此偏心球形空腔内部是匀强引力场。这个结论与静电学中均匀带电球的偏心球形空腔完全类比。

\hypertarget{ux53c2ux6570ux7075ux654fux5ea6}{%
\subsection{12.4 参数灵敏度}\label{ux53c2ux6570ux7075ux654fux5ea6}}

令 \(p=√η_1\)、\(q=√η_2\)。由

\[a/R=2q/(p+q)\]

可得一阶相对灵敏度

\[∂\ln a/∂\ln η_1=-p/[2(p+q)],\]

\[∂\ln a/∂\ln η_2=+p/[2(p+q)].\]

对半径有

\[∂\ln r/∂\ln η_1=q/[3(p+q)],\]

\[∂\ln r/∂\ln η_2=p/[3(p+q)].\]

本题 \(p>q\)，所以远点的小异常 \(η_2\)
对位置和半径反演都相当重要。实验上，越小的异常通常相对误差越大，这会成为实际反演的主要不确定度来源。

\hypertarget{ux771fux5b9eux7adeux8d5bux540cux7c7bux9898ux4e002021-fma-ux7b2c-19-ux9898}{%
\section{13. 真实竞赛同类题一：2021 F=ma 第 19
题}\label{ux771fux5b9eux7adeux8d5bux540cux7c7bux9898ux4e002021-fma-ux7b2c-19-ux9898}}

\hypertarget{ux6765ux6e90ux4e0eux5bf9ux5e94ux5173ux7cfb}{%
\subsection{13.1
来源与对应关系}\label{ux6765ux6e90ux4e0eux5bf9ux5e94ux5173ux7cfb}}

\begin{itemize}
\tightlist
\item
  赛事：2021 年美国物理队 F=ma 选拔考试。
\item
  题号：第 19 题。
\item
  官方归档：AAPT Past Exams 页面列出 2021 F=ma 试卷与官方解答。
\item
  难度：中，5/10。
\item
  核心对应：同样使用``完整球减去缺失球''；本题把观测点放在非共线位置，强迫我们进行向量分解。
\end{itemize}

\hypertarget{ux9898ux610fux91cdux8ff0}{%
\subsection{13.2 题意重述}\label{ux9898ux610fux91cdux8ff0}}

一个质量为 \(M\)、半径为 \(R\)、密度均匀的球形行星中，挖去一个半径为
\(R/2\) 的球形空腔。大球中心取 \(O=(0,0)\)，观测点 \(P=(0,R)\)
位于大球顶部；空腔中心取 \(C=(R/2,0)\)，空腔与大球在右侧内切。求点 \(P\)
处剩余天体的引力场大小。选项系数为
0.200、0.457、0.829、0.900、0.912，均乘 \(GM/R^2\)。

\hypertarget{ux5efaux6a21}{%
\subsection{13.3 建模}\label{ux5efaux6a21}}

完整大球在 \(P\) 处的场为

\[\boldsymbol g_{\text{full}}=-(GM/R^2)\boldsymbol e_y.\]

缺失球半径为 \(R/2\)，密度相同，所以质量为

\[m=M(1/2)^3=M/8.\]

从空腔中心到 \(P\) 的向量为

\[\overrightarrow{PC}=(R/2,-R),\]

距离为

\[PC=R√5/2.\]

\hypertarget{ux89e3ux6cd5ux4e00ux5206ux91cfux6cd5}{%
\subsection{13.4
解法一：分量法}\label{ux89e3ux6cd5ux4e00ux5206ux91cfux6cd5}}

缺失球原本在 \(P\) 处产生的场大小为

\[g_c=G(M/8)/(5R^2/4)=GM/(10R^2).\]

从 \(P\) 指向 \(C\) 的单位向量为

\[(1/√5,-2/√5).\]

剩余天体场等于完整球场减去缺失球原场：

\[\boldsymbol g=\boldsymbol g_{\text{full}}-\boldsymbol g_c.\]

所以

\[g_x=-(GM/R^2)/(10√5),\]

\[g_y=[-1+1/(5√5)]GM/R^2.\]

场强大小为

\[g/(GM/R^2)=√{[1/(10√5)]^2+[1-1/(5√5)]^2}.\]

数值为

\[g≈0.91165GM/R^2≈0.912GM/R^2.\]

因此选 \textbf{E}。

\hypertarget{ux89e3ux6cd5ux4e8cux4f59ux5f26ux5b9aux7406ux5f0fux7684ux5411ux91cfux6a21ux957f}{%
\subsection{13.5
解法二：余弦定理式的向量模长}\label{ux89e3ux6cd5ux4e8cux4f59ux5f26ux5b9aux7406ux5f0fux7684ux5411ux91cfux6a21ux957f}}

完整球场大小为 1 个单位，缺失球原场大小为 0.1 个单位。它们的夹角满足

\[\cos θ=2/√5.\]

剩余场是两向量之差，因此

\[g^2/(GM/R^2)^2=1^2+0.1^2-2(1)(0.1)(2/√5).\]

\[=1.01-0.4/√5.\]

开平方同样得到 0.91165，选 E。

\hypertarget{ux4e0eux4e3bux95eeux9898ux7684ux5bf9ux5e94ux5173ux7cfbux4e0eux53d8ux5f0fux70b9}{%
\subsection{13.6
与主问题的对应关系与变式点}\label{ux4e0eux4e3bux95eeux9898ux7684ux5bf9ux5e94ux5173ux7cfbux4e0eux53d8ux5f0fux70b9}}

\begin{itemize}
\tightlist
\item
  共同模型：负密度球叠加。
\item
  主问题只在共线的极值点测量，可用标量差；F=ma
  题在离轴点测量，必须做向量运算。
\item
  新卡点：减去的是``缺失球原本产生的向量场''，不是简单从完整场强大小中减去
  0.1。
\end{itemize}

\hypertarget{ux9ad8ux9891ux9519ux8bef}{%
\subsection{13.7 高频错误}\label{ux9ad8ux9891ux9519ux8bef}}

\begin{itemize}
\tightlist
\item
  把空腔质量写成 \(M/4\)，忘记体积是三次方。
\item
  把 \(PC\) 误写成 \(R\)。
\item
  直接算 \(1-0.1=0.9\)，忽略两场不共线。
\end{itemize}

\hypertarget{ux771fux5b9eux7adeux8d5bux540cux7c7bux9898ux4e8c2000-oibf-ux7406ux8bbaux7b2c-4-ux9898}{%
\section{14. 真实竞赛同类题二：2000 OIbF 理论第 4
题}\label{ux771fux5b9eux7adeux8d5bux540cux7c7bux9898ux4e8c2000-oibf-ux7406ux8bbaux7b2c-4-ux9898}}

\hypertarget{ux6765ux6e90ux4e0eux6838ux9a8cux72b6ux6001}{%
\subsection{14.1
来源与核验状态}\label{ux6765ux6e90ux4e0eux6838ux9a8cux72b6ux6001}}

\begin{itemize}
\tightlist
\item
  赛事：第五届伊比利亚美洲物理奥林匹克，2000 年，西班牙 Jaca。
\item
  题号：理论第 4 题，10 分。
\item
  西班牙语题名：\emph{Un uso práctico de la medida de la gravedad}。
\item
  官方题卷：\texttt{https://olimpicos.net/files/oibf/OIbF\_2000\_Q1-5.pdf}，合订卷第
  4---5 页。
\item
  本轮公开归档中未取得官方解答或评分细则，所以下列结果均标为独立复算，不冒充官方答案。
\item
  难度：中，5/10。
\end{itemize}

\textbf{西班牙语短节选：} ``Prospecciones geofísicas indican que
\ldots{} existe una cavidad aproximadamente esférica \ldots{}''

\hypertarget{ux5b8cux6574ux4e2dux6587ux9898ux610fux91cdux8ff0-1}{%
\subsection{14.2
完整中文题意重述}\label{ux5b8cux6574ux4e2dux6587ux9898ux610fux91cdux8ff0-1}}

地球地壳中存在一个近似球形空腔，平均半径
\(R_c=2.1\,\mathrm{km}\)。空腔内为甲烷，其密度远小于周围地球物质密度。在地表上离空腔最近的点
\(P\) 测得重力

\[g=0.99997g_0,\]

其中 \(g_0\) 是把地球视为半径 \(R_T=6370\,\mathrm{km}\) 的均匀球体时，点
\(P\) 的正常重力。忽略地球自转。

\begin{enumerate}
\def\labelenumi{\arabic{enumi}.}
\tightlist
\item
  估算空腔中心到点 \(P\) 的距离 \(d\)。
\item
  用摆长精确已知的单摆测量 \(g\)。若要求重力绝对误差不超过
  \(Δg=10^{-4}\,\mathrm{N/kg}\)，周期的相对精度 \(ΔT/T\)
  至少应达到多少？
\end{enumerate}

\hypertarget{ux4e0eux4e3bux95eeux9898ux7684ux5bf9ux5e94ux5173ux7cfb}{%
\subsection{14.3
与主问题的对应关系}\label{ux4e0eux4e3bux95eeux9898ux7684ux5bf9ux5e94ux5173ux7cfb}}

两题都把空腔视为缺失质量球。主问题给了近点和远点两个异常，因此能同时反演位置与半径；OIbF
只给一个测点，但直接给出空腔半径，所以仍可求深度。若半径也未知，单点异常只能确定
\(R_c^3/d^2\)，会出现``大小---深度退化''。

\hypertarget{ux7b2cux4e00ux95eeux89e3ux6cd5ux4e00ux5bc6ux5ea6ux6d88ux5143}{%
\subsection{14.4
第一问解法一：密度消元}\label{ux7b2cux4e00ux95eeux89e3ux6cd5ux4e00ux5bc6ux5ea6ux6d88ux5143}}

重力相对减少量为

\[ε=(g_0-g)/g_0=1-0.99997=3.0×10^{-5}.\]

设地球均匀密度为 \(ρ\)。完整地球表面重力为

\[g_0=(4πGρ/3)R_T.\]

空腔缺失质量为

\[m_c=(4π/3)ρR_c^3.\]

点 \(P\) 在空腔球外，故异常大小为

\[g_0-g=Gm_c/d^2=(4πGρ/3)R_c^3/d^2.\]

除以 \(g_0\)：

\[ε=R_c^3/(R_Td^2).\]

因此

\[d=√[R_c^3/(εR_T)].\]

代入千米单位：

\[d=√[(2.1)^3/(3.0×10^{-5}×6370)]\,\mathrm{km}.\]

\[d=6.96\,\mathrm{km}≈7.0\,\mathrm{km}.\]

空腔顶部到地表的净深度为

\[d-R_c=4.86\,\mathrm{km}.\]

\hypertarget{ux7b2cux4e00ux95eeux89e3ux6cd5ux4e8cux8d28ux91cfux6bd4ux4e0eux8dddux79bbux6bd4}{%
\subsection{14.5
第一问解法二：质量比与距离比}\label{ux7b2cux4e00ux95eeux89e3ux6cd5ux4e8cux8d28ux91cfux6bd4ux4e0eux8dddux79bbux6bd4}}

均匀密度下

\[m_c/M_T=(R_c/R_T)^3.\]

两处场强之比为

\[ε=(Gm_c/d^2)/(GM_T/R_T^2).\]

整理：

\[ε=(R_c/R_T)^3(R_T/d)^2=R_c^3/(R_Td^2).\]

于是仍得

\[d=6.96\,\mathrm{km}.\]

这条路线更短，适合考场复算；第一条路线更清楚地展示了密度和常数如何消去。

\hypertarget{ux7b2cux4e8cux95eeux89e3ux6cd5ux4e00ux5bf9ux6570ux5faeux5206}{%
\subsection{14.6
第二问解法一：对数微分}\label{ux7b2cux4e8cux95eeux89e3ux6cd5ux4e00ux5bf9ux6570ux5faeux5206}}

小振幅单摆周期为

\[T=2π√(\ell/g).\]

改写为

\[g=4π^2\ell/T^2.\]

摆长 \(\ell\) 精确已知，对数微分：

\[dg/g=-2dT/T.\]

取误差大小：

\[ΔT/T=(1/2)Δg/g.\]

取地表 \(g≈9.8\,\mathrm{N/kg}\)：

\[ΔT/T≈10^{-4}/(2×9.8)=5.10×10^{-6}.\]

即约 5.1 ppm。

\hypertarget{ux7b2cux4e8cux95eeux89e3ux6cd5ux4e8cux7cbeux786eux6270ux52a8ux68c0ux67e5}{%
\subsection{14.7
第二问解法二：精确扰动检查}\label{ux7b2cux4e8cux95eeux89e3ux6cd5ux4e8cux7cbeux786eux6270ux52a8ux68c0ux67e5}}

若周期测量值为 \(T_m=T(1+δ)\)，则

\[g_m/g=(1+δ)^{-2}.\]

展开：

\[(1+δ)^{-2}=1-2δ+3δ^2-\cdots.\]

当 \(δ≈5.1×10^{-6}\) 时，二次项约为
\(7.8×10^{-11}\)，远小于目标精度，所以一阶公式完全足够。

\hypertarget{oibf-ux9010ux5c0fux95eeux8986ux76d6ux77e9ux9635}{%
\subsection{14.8 OIbF
逐小问覆盖矩阵}\label{oibf-ux9010ux5c0fux95eeux8986ux76d6ux77e9ux9635}}

\begin{longtable}[]{@{}
  >{\raggedright\arraybackslash}p{(\columnwidth - 8\tabcolsep) * \real{0.2000}}
  >{\raggedright\arraybackslash}p{(\columnwidth - 8\tabcolsep) * \real{0.2000}}
  >{\raggedright\arraybackslash}p{(\columnwidth - 8\tabcolsep) * \real{0.2000}}
  >{\raggedright\arraybackslash}p{(\columnwidth - 8\tabcolsep) * \real{0.2000}}
  >{\raggedright\arraybackslash}p{(\columnwidth - 8\tabcolsep) * \real{0.2000}}@{}}
\toprule\noalign{}
\begin{minipage}[b]{\linewidth}\raggedright
小问
\end{minipage} & \begin{minipage}[b]{\linewidth}\raggedright
条件
\end{minipage} & \begin{minipage}[b]{\linewidth}\raggedright
定律
\end{minipage} & \begin{minipage}[b]{\linewidth}\raggedright
独立结论
\end{minipage} & \begin{minipage}[b]{\linewidth}\raggedright
完成状态
\end{minipage} \\
\midrule\noalign{}
\endhead
\bottomrule\noalign{}
\endlastfoot
1 &
\(R_c=2.1\,\mathrm{km}\)，\(R_T=6370\,\mathrm{km}\)，\(g/g_0=0.99997\) &
球壳定理、叠加、比例消元 & \(d=6.96\,\mathrm{km}\) & 完成 \\
2 & 单摆摆长精确，\(Δg=10^{-4}\,\mathrm{N/kg}\) & 单摆周期、误差传播 &
\(ΔT/T≈5.1×10^{-6}\) & 完成 \\
\end{longtable}

\hypertarget{ux53efux590dux7528ux65b9ux6cd5ux8bbaux603bux7ed3}{%
\section{15.
可复用方法论总结}\label{ux53efux590dux7528ux65b9ux6cd5ux8bbaux603bux7ed3}}

面对``均匀球 + 球形空腔 + 引力测量''题，固定执行以下流程：

\begin{enumerate}
\def\labelenumi{\arabic{enumi}.}
\tightlist
\item
  \textbf{补全母体：} 把实际物体写成完整均匀球加一个负密度球。
\item
  \textbf{先判球内还是球外：}
  球外用点质量场，球内用与位置成正比的线性场。
\item
  \textbf{画清中心距离：} 平方反比律永远使用测点到空腔中心的距离。
\item
  \textbf{优先取比值：} 多测点数据可消去
  \(G\)、密度和缺失质量，先反演几何位置。
\item
  \textbf{再用绝对异常：} 得到位置后，再求体积、半径或密度差。
\item
  \textbf{离轴必须矢量化：} 只有共线测点才可直接比较场强大小。
\item
  \textbf{反代全部观测：} 一个答案必须同时复现所有测量值。
\item
  \textbf{检查几何可行性：}
  空腔边界不能越出母球，且所用球内/球外公式必须自洽。
\end{enumerate}

\hypertarget{ux4e3bux95eeux9898ux9010ux5c0fux95eeux8986ux76d6ux77e9ux9635}{%
\section{16.
主问题逐小问覆盖矩阵}\label{ux4e3bux95eeux9898ux9010ux5c0fux95eeux8986ux76d6ux77e9ux9635}}

虽然原题没有用 a、b、c 分列，但它包含三个明确输出，必须分别覆盖。

\begin{longtable}[]{@{}
  >{\raggedright\arraybackslash}p{(\columnwidth - 10\tabcolsep) * \real{0.1667}}
  >{\raggedright\arraybackslash}p{(\columnwidth - 10\tabcolsep) * \real{0.1667}}
  >{\raggedright\arraybackslash}p{(\columnwidth - 10\tabcolsep) * \real{0.1667}}
  >{\raggedright\arraybackslash}p{(\columnwidth - 10\tabcolsep) * \real{0.1667}}
  >{\raggedright\arraybackslash}p{(\columnwidth - 10\tabcolsep) * \real{0.1667}}
  >{\raggedright\arraybackslash}p{(\columnwidth - 10\tabcolsep) * \real{0.1667}}@{}}
\toprule\noalign{}
\begin{minipage}[b]{\linewidth}\raggedright
输出要求
\end{minipage} & \begin{minipage}[b]{\linewidth}\raggedright
原题数据
\end{minipage} & \begin{minipage}[b]{\linewidth}\raggedright
关键关系
\end{minipage} & \begin{minipage}[b]{\linewidth}\raggedright
最终答案
\end{minipage} & \begin{minipage}[b]{\linewidth}\raggedright
图形依赖
\end{minipage} & \begin{minipage}[b]{\linewidth}\raggedright
状态
\end{minipage} \\
\midrule\noalign{}
\endhead
\bottomrule\noalign{}
\endlastfoot
空腔中心所在方向 & \(A\) 为最小重力点，\(B\) 为对径最大点 &
异常随到空腔中心距离减小而增大 & 中心在直径 \(AO\) 上、偏向 \(A\) & 依赖
\(A-C-O-B\) 共线图 & 完成 \\
中心埋深 \(a\) & \(η_1=0.062\)，\(η_2=0.007\) & \(√(η_1/η_2)=(2R-a)/a\)
& \(a=0.503006R≈0.503R\) & 使用 \(AC=a\)、\(BC=2R-a\) & 完成 \\
空腔半径 \(r\) & 已求 \(a\) 与任一异常值 & \((r/R)^3=η_1(a/R)^2\) &
\(r=0.250330R≈0.250R\) & 使用球外场与体积比 & 完成 \\
几何自洽 & 空腔不出露 & \(a-r>0\) & \(a-r=0.252676R\) & 近侧边界 &
完成 \\
\end{longtable}

\hypertarget{ux6700ux7ec8ux7b54ux6848ux6c47ux603b}{%
\section{17. 最终答案汇总}\label{ux6700ux7ec8ux7b54ux6848ux6c47ux603b}}

\hypertarget{ux4e3bux95eeux9898}{%
\subsection{主问题}\label{ux4e3bux95eeux9898}}

\[a/R=2√η_2/(√η_1+√η_2)=0.5030060602,\]

\[r/R=[4η_1η_2/(√η_1+√η_2)^2]^{1/3}=0.2503298898.\]

因此：

\begin{itemize}
\tightlist
\item
  空腔中心位于 \(AO\) 上，朝 \(A\) 偏离原球心；
\item
  中心埋深 \(a≈0.503R\)；
\item
  空腔半径 \(r≈0.250R\)；
\item
  中心距原球心 \(OC≈0.497R\)；
\item
  最近边界距表面约 \(0.253R\)。
\end{itemize}

官方三位小数答案与独立复算一致。官方解答中点 \(B\)
异常场的单行公式漏写密度
\(ρ\)，属量纲可识别的排版笔误，后续方程和最终答案均使用正确关系。

\hypertarget{ux540cux7c7bux9898ux4e00}{%
\subsection{同类题一}\label{ux540cux7c7bux9898ux4e00}}

2021 F=ma 第 19 题：

\[g_P≈0.912GM/R^2,\]

选 E。

\hypertarget{ux540cux7c7bux9898ux4e8c}{%
\subsection{同类题二}\label{ux540cux7c7bux9898ux4e8c}}

2000 OIbF 理论第 4 题：

\[d≈6.96\,\mathrm{km},\]

\[ΔT/T≈5.1×10^{-6}.\]

\hypertarget{ux4e09ux4e2aux5173ux952eux7a81ux7834ux53e3ux4e0eux4e09ux4e2aux9ad8ux9891ux9519ux8bef}{%
\section{18.
三个关键突破口与三个高频错误}\label{ux4e09ux4e2aux5173ux952eux7a81ux7834ux53e3ux4e0eux4e09ux4e2aux9ad8ux9891ux9519ux8bef}}

\hypertarget{ux4e09ux4e2aux5173ux952eux7a81ux7834ux53e3}{%
\subsection{三个关键突破口}\label{ux4e09ux4e2aux5173ux952eux7a81ux7834ux53e3}}

\begin{enumerate}
\def\labelenumi{\arabic{enumi}.}
\tightlist
\item
  有空腔天体写成``完整均匀球减去缺失小球''。
\item
  用近、远两点异常之比先消去空腔大小，反演中心位置。
\item
  区分共线标量问题与离轴向量问题。
\end{enumerate}

\hypertarget{ux4e09ux4e2aux6700ux9ad8ux9891ux9519ux8bef}{%
\subsection{三个最高频错误}\label{ux4e09ux4e2aux6700ux9ad8ux9891ux9519ux8bef}}

\begin{enumerate}
\def\labelenumi{\arabic{enumi}.}
\tightlist
\item
  把中心距离 \(a\) 错写为边界距离 \(a-r\)。
\item
  把质量比例写成 \((r/R)^2\)，或忘记球体体积的三次方。
\item
  把近似 \(a≈R/2\)、\(r≈R/4\) 写成精确等号。
\end{enumerate}

\hypertarget{ux6765ux6e90ux4e0eux53c2ux8003ux8d44ux6599}{%
\section{19.
来源与参考资料}\label{ux6765ux6e90ux4e0eux53c2ux8003ux8d44ux6599}}

\begin{enumerate}
\def\labelenumi{\arabic{enumi}.}
\tightlist
\item
  BelPhO 2003
  年共和国物理奥林匹克年度归档：\texttt{https://belpho.org/problems/final/2003/}。
\item
  BelPhO 2003 年 10 年级第 2
  题独立页：\texttt{https://belpho.org/problems/final/2003/10-2/}。
\item
  官方题卷：\texttt{https://belpho.org/problems/final/2003problems.pdf}，第
  3 页。
\item
  官方解答：\texttt{https://belpho.org/problems/final/2003solutions.pdf}，第
  6 页。
\item
  AAPT Past
  Exams：\texttt{https://www.aapt.org/Common2022/pastexams.cfm}，2021
  F=ma 试卷与解答。
\item
  OIbF 归档：\texttt{https://olimpicos.net/oibf/}。
\item
  OIbF 2000
  理论题合订卷：\texttt{https://olimpicos.net/files/oibf/OIbF\_2000\_Q1-5.pdf}。
\end{enumerate}

\hypertarget{ux5f53ux524dux5b8cux6210ux72b6ux6001}{%
\section{20. 当前完成状态}\label{ux5f53ux524dux5b8cux6210ux72b6ux6001}}

\begin{itemize}
\tightlist
\item[$\boxtimes$]
  数据库序号与官方题号映射核验。
\item[$\boxtimes$]
  官方题面位置与官方解答位置核验。
\item[$\boxtimes$]
  完整中文题意重述与符号说明。
\item[$\boxtimes$]
  读题拆解、六步建模、突破口分析。
\item[$\boxtimes$]
  两种组织不同的完整解法。
\item[$\boxtimes$]
  官方答案与独立复算对照。
\item[$\boxtimes$]
  官方局部排版笔误说明。
\item[$\boxtimes$]
  量纲、反代、极限、几何和物理合理性检查。
\item[$\boxtimes$]
  常见误区与正确改法。
\item[$\boxtimes$]
  两道真实同类竞赛题的来源、对应关系和完整推导。
\item[$\boxtimes$]
  一般化、参数可辨识性、向量场、空腔内匀强场与灵敏度分析。
\item[$\boxtimes$]
  逐输出覆盖矩阵与最终答案汇总。
\end{itemize}

\textbf{本文应命名为：}
\texttt{2003\_白俄罗斯共和国物理奥林匹克决赛\_10年级第2题\_小行星球形空腔反演\_完整精讲.md}

\end{document}
